% ============================================================
%  Polynomial Illumination Correction (Flat Field)
%  for the Arsenic Colorimeter
%
%  Companion to colorimeter_methods.tex (Layer 1) and
%  colorimeter_plan.tex.  Documents the analysis carried out with
%  devdesk/flatfield.py on the 2026-08-03 uniform-card suites.
%
%  Source formatting: each sentence, and each major clause within a
%  long sentence, begins on a new source line.
%
%  Compile with pdflatex (twice for cross-references).
% ============================================================

\documentclass[11pt, a4paper]{article}

\usepackage{amsmath, amssymb}
\usepackage{booktabs}
\usepackage{geometry}
\geometry{margin=2.5cm}
\usepackage{hyperref}
\hypersetup{colorlinks=true, linkcolor=blue, citecolor=blue, urlcolor=blue}
\usepackage{microtype}
\usepackage{parskip}
\usepackage{enumitem}

\newcommand{\code}[1]{\texttt{#1}}
\newcommand{\pct}{\,\%}
\newcommand{\mean}[2]{\left\langle #1 \right\rangle_{#2}}

\title{%
  Polynomial Illumination Correction for the Arsenic Colorimeter:\\
  Analysis, Criteria, Implementation and Protocol}
\author{Development Notes}
\date{2026-08-04}

\begin{document}
\maketitle
\tableofcontents
\newpage

% ============================================================
\section{Purpose and scope}
% ============================================================

The colorimeter reads a chemical test strip by comparing the mean
brightness of a small \emph{sample} patch
against the mean brightness of a \emph{white reference} patch
elsewhere on the same card.
That comparison is meaningful only if the two patches are lit equally.
They are not.
Measured on the captures described below,
the illumination inside the imaged region falls by about $25\pct$
from the top of the region to the bottom.\footnote{Denominator:
the mean over the whole imaged region.
The figure is the difference between the mean of the top row of pixels
and the mean of the bottom row, divided by that overall mean.}

This document covers the correction for that non-uniformity.
The correction map is called a \emph{flat
field}\footnote{\textbf{Flat field.}
A map of how brightly the illumination lands on each pixel
when the imaged object is uniform.
Dividing a picture by its flat field removes the lighting pattern
and leaves only the object.
The name is from astronomy,
where the calibration image is taken of a featureless (``flat'') screen.},
and the decision documented here is to represent it as a low-order
polynomial surface,
fitted separately for each LED color and each Bayer layer.

The document reports:

\begin{itemize}[nosep]
  \item the notation and the definition of every measure used
        (\S\ref{sec:notation});
  \item what the correction is and where it sits in the pipeline
        (\S\ref{sec:pipeline});
  \item whether the correction has to be a map at all,
        or whether two numbers would do (\S\ref{sec:scalar});
  \item the polynomial model and why it takes the form it does
        (\S\ref{sec:model});
  \item the criteria proposed for judging a fit adequate
        (\S\ref{sec:criteria});
  \item what was measured, on which captures (\S\ref{sec:results});
  \item the choices those measurements drove (\S\ref{sec:choices});
  \item how the evaluation harness is built (\S\ref{sec:implementation});
  \item the image sampling protocol for the captures still needed
        (\S\ref{sec:protocol});
  \item how to run the programs and read their output
        (\S\ref{sec:running}, \S\ref{sec:interpreting});
  \item what this unblocks (\S\ref{sec:next}).
\end{itemize}

Nothing described here has been wired into the instrument's capture or
analysis path.
The work so far is a desktop evaluation harness,
\code{devdesk/flatfield.py},
plus the archive metadata needed to make its inputs self-describing.

% ============================================================
\section{Notation, vocabulary and measures}
\label{sec:notation}
% ============================================================

\subsection{Terms}

\begin{description}[style=nextline, leftmargin=1em]

\item[Bayer layer]
  The camera sensor has one color filter per photosite,
  in a repeating $2\times2$ tile:
  red, two greens, blue.
  Splitting the raw image on that tile gives four half-resolution images,
  called \emph{layers} and labeled \code{R}, \code{Gr}, \code{Gb}, \code{B}.
  \code{Gr} is the green photosite that shares a row with red;
  \code{Gb} is the one that shares a row with blue.
  They are not interchangeable;
  see \S\ref{sec:crosstalk}.

\item[DN]
  Digital Number: the raw sensor count, before conversion to physical units.
  The sensor is 10-bit, so DN runs $0$ to $1023$.

\item[Stream space]
  Pixel coordinates in the full camera frame,
  $2304\times1296$ for the IMX708 mode in use.

\item[Layer space]
  Pixel coordinates within one Bayer layer.
  A rectangle $w\times h$ in stream space is $w/2 \times h/2$ in layer space.
  All analysis in this document is in layer space.

\item[\code{bound}]
  The rectangular region of interest that the capture path crops and stores.
  Everything outside it is discarded at capture time.
  Currently $196\times660$ in stream space,
  which is $98\times330$ per layer.

\item[Crop]
  The stored \code{bound} region of one Bayer layer,
  as an array of $330$ rows by $98$ columns.
  ``The crop'' always means this whole array.

\item[Patch]
  One of the small rectangles the analysis measures:
  \code{sample}, \code{white}, \code{chip1}, \code{chip2}.
  Each is $68\times68$ in stream space,
  so $34\times34$ in layer space, which is $1156$ pixels.
  A patch is a real region of interest at a real location on the card.
  Patches are addressed in code by \emph{role}
  (\code{ROLE\_SAMPLE}, \code{ROLE\_WHITE}, \ldots), never by name,
  because names are operator-editable.

\item[Window]
  One tile of a regular $34\times34$ grid laid over the crop,
  used only as a measuring grid for the block-residual statistic
  (\S\ref{sec:m-block}).
  Windows are the same size as patches
  but are placed by the grid, not by where an ROI is.
  The $330\times98$ crop holds $9\times2 = 18$ whole windows.

\item[Extent]
  The set of pixels a polynomial was fitted over,
  written $\Omega$.
  The extent may be the whole crop or some subset of it;
  the choices tried are listed in \S\ref{sec:ladder}.

\item[Dark frame]
  A capture taken at the same exposure with the LEDs off.
  Subtracting it removes the sensor's electronic offset and thermal signal.
  All work here is on dark-subtracted data.

\item[Suite]
  One capture archive (\code{.h5}) holding,
  for every LED color,
  a ladder of exposures plus the matching dark frames.

\item[In-sample]
  A fit evaluated on the same capture it was fitted to.
  Optimistic by construction:
  the fit has already seen that capture's noise and texture.

\item[Out-of-sample, or transfer]
  A fit taken from one capture and applied to a different one.
  This is what the instrument does in service,
  because it applies a stored calibration to a new measurement.
  It is therefore the estimate that predicts field performance.

\end{description}

\subsection{Symbols}
\label{sec:symbols}

Every quantity below is defined per Bayer layer,
for one LED color at one exposure.
There is no mixing of layers or colors anywhere in the analysis.

\begin{center}
\begin{tabular}{ll}
\toprule
Symbol & Meaning \\
\midrule
$p=(x,y)$ & a pixel location in layer space; $x$ is the column, $y$ the row \\
$D(p)$    & lit frame, DN \\
$D_0(p)$  & dark frame at the same exposure, DN \\
$S(p) = D(p)-D_0(p)$ & dark-subtracted signal, DN \\
$A(p)$    & reflectance of the imaged target, dimensionless --- the quantity wanted \\
$L(p)$    & illumination delivered to $p$, DN per unit reflectance \\
$\Omega$  & the fit extent: a set of pixels \\
$|\Omega|$& the number of pixels in that set \\
$F(p)$    & the fitted polynomial surface, DN \\
$G(p)$    & the gain map, dimensionless \\
$S^{c}(p)$& corrected signal, DN \\
$R_i$     & the set of pixels in patch $i$; $n_i = |R_i|$, and $n_i = 1156$ \\
$w_i,h_i$ & the width and height of patch $i$ in layer pixels; both $34$ \\
$\bar{s}_i$ & the corrected mean of patch $i$, DN \\
$\bar{g}_i$ & the mean of the gain map over patch $i$, dimensionless \\
$W_k$     & the $k$-th window of the measuring grid \\
\bottomrule
\end{tabular}
\end{center}

The angle-bracket notation means an arithmetic mean over a set of pixels:
\begin{equation}
  \mean{X}{\Omega} \;\equiv\; \frac{1}{|\Omega|}\sum_{p \in \Omega} X(p) .
  \label{eq:mean}
\end{equation}
The subscript names the set.
$\mean{X}{\Omega}$ is a scalar.

The model is a pixelwise one throughout.
$D$, $D_0$, $S$, $A$, $L$, $F$, $G$ and $S^{c}$ are all arrays
of the same shape as the crop,
and every operation combining them --- subtraction in
Eq.~\ref{eq:sig}, division in Eq.~\ref{eq:correct} ---
is elementwise, one pixel at a time.
No matrix product appears anywhere in the correction.
The only scalars are the means defined by Eq.~\ref{eq:mean}
and the polynomial coefficients.

\subsection{The three measures}
\label{sec:measures}

Three numbers recur throughout the results.
All three are relative quantities expressed as percentages,
and the denominator of each is part of its definition.

\subsubsection{Block residual --- how flat the corrected crop is}
\label{sec:m-block}

Lay the $34\times34$ measuring grid over the corrected crop
and take the mean of each window:
\begin{equation}
  m_k \;=\; \mean{S^{c}}{W_k} .
\end{equation}
If the correction were perfect and the imaged card uniform,
every $m_k$ would equal the mean of the whole corrected crop.
Express each window's departure from that as a percentage:
\begin{equation}
  d_k \;=\; 100 \left( 
    \frac{m_k - \mean{S^{c}}{\mathrm{crop}}}{\mean{S^{c}}{\mathrm{crop}}}  
  \right) .
  \label{eq:block}
\end{equation}
The ratio $m_k / \mean{S^{c}}{\mathrm{crop}}$ equals $1$ when the window
matches the whole crop,
so subtracting $1$ converts it from a ratio into a deviation,
and multiplying by $100$ expresses that deviation as a percentage.
The denominator is the mean of the whole corrected crop.

Two summaries of $\{d_k\}$ are reported:
the root-mean-square over the $18$ windows,
printed as \code{blk rms},
and the largest $|d_k|$, printed as \code{blk max}.

Averaging over $1156$ pixels reduces the shot-noise contribution
to about $0.04\pct$ of the window mean,
so what survives in $d_k$ is spatial structure,
not pixel noise.

\subsubsection{Criterion P --- patch-ratio accuracy}
\label{sec:m-p}

$\bar{s}_i$ is the corrected mean of patch $i$,
that is $\mean{S^{c}}{R_i}$,
a single number in DN.
$\bar{s}_{\mathrm{white}}$ is the same quantity for the patch
in the white-reference role,
and it is the denominator.

Image a card that is genuinely uniform.
Every patch should then read the same as the white reference,
so for each patch $i$ the quantity
\begin{equation}
  P_i \;=\; 100 \left| 
  %\frac{\bar{s}_i}{\bar{s}_{\mathrm{white}}} - 1 \right|
  \frac{\bar{s}_i - \bar{s}_{\mathrm{white}}}{\bar{s}_{\mathrm{white}}} \right|
  \label{eq:critP}
\end{equation}
should be zero, and whatever it is instead is the flat field's error.

Two levels of aggregation follow, and they are easy to confuse:

\begin{enumerate}[nosep]
  \item Within one frame, take the worst patch:
        $P = \max_i P_i$, over the patches present.
  \item Across the set of frames scored --- one frame per
        (LED color $\times$ Bayer layer) --- report the median of $P$
        as \code{P med} and the maximum of $P$ as \code{P max}.
\end{enumerate}

So \code{P med} $= 0.24\pct$ means:
in half the 39 frames scored,
the worst-deviating patch deviated by less than $0.24\pct$
of that frame's corrected white-patch mean,
and in half of them by more.
\code{P max} is the single worst patch in the single worst frame.
Neither figure is a median over patches.

\subsubsection{Criterion S --- residual tilt inside a patch}
\label{sec:m-s}

Fit a plane to the corrected patch by least squares,
\begin{equation}
  \hat{v}(x,y) \;=\; a_0 + a_x\,x + a_y\,y ,
  \qquad (x,y) \in R_i ,
  \label{eq:plane}
\end{equation}
where $x$ runs $0 \ldots w_i-1$ and $y$ runs $0 \ldots h_i-1$
in layer pixels within the patch.
The constant $a_0$ carries the patch level;
$a_x$ and $a_y$ are the slopes, in DN per layer pixel.

A plane over a rectangle attains its largest and smallest values
at opposite corners,
so the peak-to-peak range of $\hat{v}$ over the patch is
\begin{equation}
  R_{\mathrm{plane}} \;=\; |a_x|\,(w_i - 1) \;+\; |a_y|\,(h_i - 1) ,
  \label{eq:range}
\end{equation}
in DN.
The spans are $w_i - 1$ and $h_i - 1$, not $w_i$ and $h_i$,
because a row of $w_i$ pixels spans $w_i - 1$ intervals
between first and last pixel center.
Using $w_i$ would overstate $R_{\mathrm{plane}}$
by a factor $w_i/(w_i-1)$, here $34/33$.

Criterion S is that range as a percentage of the patch's own mean:
\begin{equation}
  S_i \;=\; 100\,\frac{R_{\mathrm{plane}}}{\bar{s}_i} .
  \label{eq:critS}
\end{equation}
The denominator is patch $i$'s own corrected mean, not the white patch's.

$S_i$ needs a companion figure,
because a plane fitted to pure noise still has non-zero slopes.
Take $n$ pixels evenly spaced along a side of length $\ell$,
with independent noise of standard deviation $\sigma_v$ on each.
The least-squares slope has standard error
$\sigma_v \sqrt{12 / (n\,\ell^{2})}$,
so the slope multiplied by the side length has standard error
$\sigma_v \sqrt{12/n}$.
The two axes contribute one such term each, and Eq.~\ref{eq:range} adds them,
so the range expected from noise alone is
\begin{equation}
  S^{\mathrm{nz}}_i \;=\; 2\,\sigma_i\,\sqrt{12/n_i} ,
  \label{eq:snz}
\end{equation}
where $\sigma_i$ is the pixel standard deviation within patch $i$
as a percentage of $\bar{s}_i$.
For $n_i = 1156$ the factor $2\sqrt{12/n_i}$ is $0.204$.

The reported column \code{S/nz} is the ratio $S_i / S^{\mathrm{nz}}_i$.
A value near $1$ means the tilt is what noise alone would have produced,
so it is a fitting artifact.
A value well above $1$ means the corrected patch really does carry a tilt.

\subsection{Other percentages}

Percentages that do not fit one of the three definitions above
are spelled out where they appear.
The recurring ones are the masked-pixel fraction
(denominator: pixels in the fit extent)
and brightness ratios between two cards
(denominator: named at the point of use).

% ============================================================
\section{Where the correction sits}
\label{sec:pipeline}
% ============================================================

The pipeline for one measurement, per Bayer layer:

\begin{enumerate}[nosep]
  \item Capture the lit frame $D$ and its dark frame $D_0$
        at the same exposure.
  \item Subtract, pixel by pixel:
        \begin{equation}
          S(p) \;=\; D(p) - D_0(p) .
          \label{eq:sig}
        \end{equation}
  \item Mask pixels that are saturated or dead (\S\ref{sec:masking}).
  \item Divide $S$ by the gain map $G$ --- this document's subject.
  \item Take the patch means $\bar{s}_{\mathrm{sample}}$
        and $\bar{s}_{\mathrm{white}}$,
        and form their ratio.
  \item Convert that ratio to concentration by the strip calibration.
\end{enumerate}

Step 4 is the only step that knows anything about illumination geometry.
Its sole purpose is to make step 5's ratio independent of
\emph{where on the card} the two patches sit.

\subsection{The correction divides; it does not subtract}
\label{sec:divide}

\code{colorimeter\_methods.tex} specified subtracting a saved trend array.
That is correct only to first order.
Illumination acts multiplicatively on the dark-subtracted signal:
a pixel that receives $10\pct$ less light returns $10\pct$ less signal,
not a fixed number of DN less.
Writing the signal as the product of what the target reflects
and what the lamp delivers,
\begin{equation}
  S(p) \;=\; A(p)\,L(p) ,
  \label{eq:model}
\end{equation}
makes the point plain:
$L$ multiplies, so removing it requires division.

The size of the error from getting this wrong was measured.
On the two 2026-08-03 captures,
the gray card returned $0.27\times$ the signal of the white card
at the same exposure.\footnote{Ratio of dark-subtracted patch means,
gray divided by white.
Denominator: the white card's patch mean.}
An additive trend calibrated on the white card,
applied to the gray card,
over-corrects by a factor of $1/0.27 \approx 3.7$.

The correction used here is therefore
\begin{equation}
  S^{c}(p)
  \;=\;
  \frac{S(p)}{G(p)},
  \qquad
  G(p) \;=\; \frac{F(p)}{\mean{F}{\Omega}} .
  \label{eq:correct}
\end{equation}
Both operations are elementwise over pixels.
$F$ is the fitted polynomial surface in DN,
$\mean{F}{\Omega}$ is a single number,
and $G$ is dimensionless with mean exactly $1$ over $\Omega$.
$G$ is called a \emph{gain map}.

Note which extent normalizes $F$.
The denominator $\mean{F}{\Omega}$ is the mean of the fitted surface
over the extent it was fitted on,
not over the whole crop.
For a fit taken over the whole crop the two are the same.
For a fit taken over a small window around one patch they are not,
and using the wrong one rescales the result without any visible symptom.
The code stores this number as \code{FlatField.norm\_mean}.

% ============================================================
\section{Does the correction need to be a map?}
\label{sec:scalar}
% ============================================================

The measurement of interest is a ratio of two patch means,
and a patch mean is a single number.
It is worth asking directly whether the whole gain \emph{map} is needed,
or whether two scalars would give the same answer.
The question turns out to have a clean answer,
and answering it also explains two results reported later.

\subsection{The exact case}

Suppose the fit is perfect,
so that $G = L / \mean{L}{\Omega}$ exactly.
Substituting Eq.~\ref{eq:model} into Eq.~\ref{eq:correct},
\begin{equation}
  S^{c}(p) \;=\; \frac{A(p)\,L(p)}{L(p)/\mean{L}{\Omega}}
             \;=\; A(p)\,\mean{L}{\Omega} .
\end{equation}
The illumination cancels pixel by pixel,
leaving reflectance times a constant.
Taking patch means and forming the ratio,
\begin{equation}
  \frac{\bar{s}_{\mathrm{sample}}}{\bar{s}_{\mathrm{white}}}
  \;=\;
  \frac{\mean{A}{R_{\mathrm{sample}}}}{\mean{A}{R_{\mathrm{white}}}} ,
\end{equation}
which is exactly the reflectance ratio wanted,
and the constant $\mean{L}{\Omega}$ has cancelled.

\subsection{What happens with two scalars instead}

Now suppose the map is never applied,
and the patch means are taken from the uncorrected signal
and divided by the patch-average gains
\begin{equation}
  \bar{g}_i \;=\; \mean{G}{R_i} ,
\end{equation}
which are two numbers.
The estimate becomes
\begin{equation}
  r_{\mathrm{scalar}}
  \;=\;
  \frac{\mean{S}{R_{\mathrm{sample}}} \big/ \bar{g}_{\mathrm{sample}}}
       {\mean{S}{R_{\mathrm{white}}}  \big/ \bar{g}_{\mathrm{white}}} .
  \label{eq:rscalar}
\end{equation}

The only step that can go wrong is
$\mean{A L}{R_i} \ne \mean{A}{R_i}\mean{L}{R_i}$.
The mean of a product is the product of the means plus their covariance,
so the relative error of that step is
\begin{equation}
  \frac{\mean{AL}{R_i}}{\mean{A}{R_i}\mean{L}{R_i}} - 1
  \;=\;
  \rho_i \; c^{A}_i \; c^{L}_i ,
  \label{eq:cov}
\end{equation}
where $c^{A}_i$ and $c^{L}_i$ are the coefficients of
variation\footnote{\textbf{Coefficient of variation.}
Standard deviation divided by mean, a dimensionless measure of relative
spread.
Written as a percentage here.}
of reflectance and of illumination \emph{within} patch $i$,
and $\rho_i$ is their correlation coefficient, between $-1$ and $+1$.

So the two-scalar form is exact whenever the target reflectance
is uniform within each patch,
because then $c^{A}_i = 0$.
Where reflectance is not uniform,
the error is the product of three factors,
each of which is small.

\subsection{How large is that error here?}

$c^{L}_i$ can be measured directly:
evaluate the fitted gain map over each patch rectangle
and take its coefficient of variation.
Median over the 39 white-suite frames, \code{Gr} layer, $(3,1)$ fit:

\begin{center}
\begin{tabular}{lrrr}
\toprule
Patch & $c^{L}_i$ ($\pct$ of patch mean) & peak-to-peak ($\pct$)
      & $\bar{g}_i$ range over LEDs \\
\midrule
\code{sample} & 0.31 & 1.33 & $1.066 \ldots 1.139$ \\
\code{chip1}  & 0.82 & 3.43 & $1.042 \ldots 1.104$ \\
\code{white}  & 1.48 & 5.75 & $0.957 \ldots 0.985$ \\
\code{chip2}  & 2.47 & 9.02 & $0.706 \ldots 0.854$ \\
\bottomrule
\end{tabular}
\end{center}

Inside the sample patch the gain varies by only $0.31\pct$.
Take a pessimistic case:
a test strip whose color development varies by $20\pct$ across the patch
($c^{A} = 20\pct$),
correlated perfectly with the illumination gradient ($\rho = 1$).
Eq.~\ref{eq:cov} then gives $1.0 \times 0.20 \times 0.0031 = 0.062\pct$,
which is below the $0.12\pct$ repeatability floor of \S\ref{sec:floor}.
For the white reference patch,
where $c^{L} = 1.48\pct$ but the target is a reference surface
and $c^{A}$ should be near $1\pct$,
the same formula gives about $0.015\pct$.

Two useful observations fall out of that table.
The sample patch happens to sit where the illumination is flattest,
so the approximation is safest exactly where it matters most.
And \code{chip2}, at the bottom of \code{bound},
has eight times the sample patch's within-patch gain variation
and the largest spread of $\bar{g}$ across LED colors ---
one more reason it is the worst patch in every test reported later.

\subsection{What still needs the map}

The two-scalar form is adequate for the concentration measurement itself,
but it does not replace the map everywhere.

\begin{itemize}[nosep]
  \item \textbf{Criterion S and the Layer 3 uniformity flag.}
        Both are statements about variation \emph{within} a patch,
        which a patch-average scalar has by definition discarded.
  \item \textbf{Masked or robust patch statistics.}
        If some pixels in a patch are excluded ---
        hot pixels, a bubble, a defect ---
        the correct scalar is the gain averaged over the surviving pixels,
        not over the whole rectangle.
        Computing that requires the map.
        The same applies to a median or trimmed mean,
        for which the cancellation argument above does not hold,
        because those estimators are not linear.
  \item \textbf{Any diagnostic that looks at the corrected image.}
        The residual heatmaps and profiles of \S\ref{sec:running} are examples.
\end{itemize}

Note also what the two-scalar form does \emph{not} avoid:
the polynomial itself.
$\bar{g}_i$ is the average of the fitted surface over patch $i$,
and there is no way to obtain it without fitting a surface.
What the argument removes is the per-pixel division at measurement time,
not the calibration.

\subsection{Has this been established?}

Partly, and it is worth being exact about which part.

\textbf{Established.}
Criterion P, Eq.~\ref{eq:critP}, is defined on patch means.
It therefore already measures the accuracy of precisely the scalar ratio
that Eq.~\ref{eq:rscalar} would use.
Every number reported in \S\ref{sec:results} is a statement about
that ratio,
not about pixelwise flatness.
The criterion was designed that way before this question was asked,
and the question confirms the design rather than changing it.

\textbf{Established.}
The within-patch gain variation $c^{L}_i$ is measured, and it is small
for the sample patch.

\textbf{Not established.}
The within-patch reflectance variation $c^{A}$ of a developed test strip,
and its correlation $\rho$ with the illumination gradient,
have not been measured.
A uniform card cannot measure them,
because a uniform card has $c^{A} \approx 0$ by assumption.
That measurement requires imaging developed strips,
which is a separate exercise.
The bound of $0.062\pct$ above is an estimate under an assumed
$c^{A} = 20\pct$, not a measurement.

\textbf{Consequence for the design.}
Store, per (LED color, Bayer layer, patch role), the single scalar
$\bar{g}_i$;
that is what the instrument needs at measurement time,
and it is a few dozen numbers rather than a few dozen arrays.
Keep the polynomial coefficients as the source those scalars are computed
from, so that a moved ROI is a recomputation rather than a recapture.
Keep the map generation available for diagnostics and for criterion S.

\subsection{Why the ratio form explains the multi-field failure}
\label{sec:whyratio}

Eq.~\ref{eq:rscalar} depends on the two gains only through their
\emph{ratio} $\bar{g}_{\mathrm{sample}}/\bar{g}_{\mathrm{white}}$.
Any factor common to both cancels.
That is why the choice of $\mean{F}{\Omega}$ in Eq.~\ref{eq:correct}
does not affect the answer
when a single field corrects both patches:
it is the same factor top and bottom.

It also predicts the failure reported in \S\ref{sec:localfail}.
When the sample patch and the white patch are corrected by
\emph{different} fields,
fitted over different extents,
the two normalizing constants $\mean{F}{\Omega_{\mathrm{sample}}}$
and $\mean{F}{\Omega_{\mathrm{white}}}$ are different numbers
and do not cancel.
Their ratio is exactly the relative brightness of the two extents,
which is the quantity being measured.
A multi-field scheme therefore cannot produce a patch ratio at all
without a further step that ties the fields to a common scale.

% ============================================================
\section{The model}
\label{sec:model}
% ============================================================

\subsection{Polynomial form}

For one combination of (LED color, exposure, Bayer layer),
the fitted surface is
\begin{equation}
  F(x,y) \;=\;
  \sum_{i_y=0}^{d_y} \; \sum_{i_x=0}^{d_x}
  c_{i_y i_x}\, x^{i_x}\, y^{i_y},
  \qquad
  \text{subject to } i_x + i_y \le \max(d_y, d_x),
  \label{eq:poly}
\end{equation}
where $d_y$ and $d_x$ are the maximum exponents allowed on the two axes,
$c_{i_y i_x}$ are the coefficients found by least squares,
and $x$ and $y$ are pixel coordinates rescaled linearly to the interval
$[-1,1]$.

The rescaling is always taken over the whole \code{bound} crop,
whatever extent the fit was taken over.
That keeps coefficient vectors from different extents directly comparable.
Fitting is ordinary least squares over the unmasked pixels of $\Omega$.

The constraint $i_x + i_y \le \max(d_y,d_x)$,
applied on top of the separate limits $d_y$ and $d_x$,
is what makes the basis \emph{anisotropic}.\footnote{\textbf{Anisotropic
basis.}
A basis that allows different maximum powers on the two axes.
A conventional bivariate polynomial of degree $n$ is isotropic:
it treats $x$ and $y$ alike and costs $\binom{n+2}{2}$ terms.
Here the surface varies strongly in $y$ and barely in $x$,
so high powers of $x$ add parameters without adding descriptive power.}
With $(d_y,d_x) = (3,1)$ it admits 7 terms,
\[
  1,\quad x,\quad y,\quad xy,\quad y^{2},\quad xy^{2},\quad y^{3},
\]
against 10 for a full bivariate cubic and 15 for a full quartic.

\subsection{Why $(3,1)$}

The surface is close to one-dimensional.
Averaging each row of the crop and comparing top to bottom
gives a fall of about $25\pct$
(denominator: the crop mean).
Averaging each column and comparing left to right
gives a variation of about $2\pct$
(same denominator).
The $y$ axis carries more than ten times the variation of the $x$ axis.

The measured consequence, on the white suite:
\code{P med} improves from $0.24\pct$ at $(3,1)$ to $0.21\pct$ at $(3,3)$.
That is a gain of $0.03\pct$ of the corrected white-patch mean,
against a tolerance of $0.5\pct$,
in exchange for three more free parameters.
See Table~\ref{tab:ladder-white}.

\subsection{One polynomial per LED color and per Bayer layer}

Each LED \emph{color} is a symmetrical array of four emitters.
Arrays of different colors cannot occupy the same physical positions,
so each color lights the stage from a slightly different geometry,
in addition to part-to-part differences between individual emitters.
A separate polynomial per color is therefore expected by construction.

The data agrees.
Comparing two normalized fields as $34\times34$ block means of their ratio
(denominator: the frame mean),
field shapes differ between LED colors,
on the same card,
by $0.55$ to $2.1\pct$.
On the same measure,
\code{R} and \code{Gr} field shapes on the \emph{same frame}
differ by up to $0.96\pct$,
so layers must be separated as well.

\subsection{Masking}
\label{sec:masking}

Three classes of pixel are excluded before fitting:

\begin{itemize}[nosep]
  \item \textbf{Saturated}: raw value $\ge 1000$ DN.
        The threshold is $1023 - 23$.
        The top of the range is already in the sensor's non-linear region,
        so those counts understate the light that produced them.
  \item \textbf{Non-positive after dark subtraction}: no signal to fit.
  \item \textbf{Hot pixels}, as flagged by \code{devdesk/roi\_darkdiff.py}.
\end{itemize}

A frame with more than $1\pct$ of its pixels masked
(denominator: pixels in the fit extent)
is refused rather than fitted,
because the surviving pixels are systematically the dimmer ones
and a fit over them is biased.

Masking changes conclusions here; it is not routine housekeeping.
The stored exposure ladders were metered against a \emph{sample}
at the center of the frame,
so a white card filling the whole stage oversaturates.
In the white suite,
\code{royal\_blue/B} has a mean of 959\,DN and,
without masking,
fits with a residual of $0.00\pct$ ---
a perfect fit to a clipped plateau.
\code{green/e2} in \code{Gr} and \code{Gb} is $96\pct$ masked
(denominator: pixels in the crop)
and is correctly refused;
see the \code{SKIPPED} rows in \S\ref{sec:running}.

% ============================================================
\section{Criteria for adequacy}
\label{sec:criteria}
% ============================================================

The flat field's only job is to make patch \emph{means} comparable
between locations.
Two consequences shape the criteria,
and \S\ref{sec:scalar} has already made both precise:

\begin{itemize}[nosep]
  \item An error that is constant over a patch cancels in the
        sample-to-white ratio,
        because it appears in both numerator and denominator.
  \item An error that is a linear ramp across a patch does not bias the
        patch mean,
        because a plane's departures from its center value are
        antisymmetric about that center and cancel in the average.
\end{itemize}

So pixelwise flatness is the wrong basis for a criterion.
Judged pixelwise, models that are entirely adequate would fail,
and overfitted models would score well.
The criteria are defined on patch means instead.

\subsection{P --- patch-ratio accuracy (primary)}

Defined by Eq.~\ref{eq:critP} in \S\ref{sec:m-p}.
Require $P \le \mathrm{tol}_P$,
with two tolerances, because they answer different questions:

\begin{center}
\begin{tabular}{lll}
\toprule
Mode & $\mathrm{tol}_P$ & Rationale \\
\midrule
in-sample & $0.5\pct$ & about $4\times$ the $0.12\pct$ repeatability floor \\
transfer  & $1.0\pct$ & what a stored calibration must achieve in service \\
\bottomrule
\end{tabular}
\end{center}

Both are percentages of the corrected white-reference patch mean
$\bar{s}_{\mathrm{white}}$,
which is the same denominator the instrument's own answer uses.
P is the criterion that decides acceptance.

\subsection{S --- residual tilt inside a patch (secondary)}

Defined by Eq.~\ref{eq:critS} in \S\ref{sec:m-s},
with the noise reference of Eq.~\ref{eq:snz}.
Require $S_i \le 1\pct$ of patch $i$'s own corrected mean.

As established in \S\ref{sec:scalar},
a tilt does not bias the patch mean,
so S does not affect the concentration measurement.
It matters because Layer 3 of \code{colorimeter\_plan.tex}
uses within-patch spread as a strip-uniformity flag,
and a leftover illumination tilt would inflate that flag
into false alarms.

\subsection{T --- degree selection (tertiary)}

Accept the smallest $(d_y,d_x)$ whose \emph{transfer} P
stops improving by more than $0.1\pct$,
that figure being the measured repeatability floor of \S\ref{sec:floor}.

The purpose of T is to prevent one specific mistake.
In-sample residual falls monotonically as degree rises, without limit,
because a more flexible model can always follow one capture more closely.
Transfer error stops falling once the added terms are describing
that capture's own noise and texture rather than the illumination.
Selecting on in-sample residual therefore guarantees an overfitted model.
The harness prints in-sample and transfer side by side in every table
so the difference between them stays visible.

% ============================================================
\section{The extent ladder}
\label{sec:ladder}
% ============================================================

``Fit over an area larger than the target'' leaves open how much larger.
Five rungs were defined, cheapest first,
to be climbed only as far as necessary:

\begin{enumerate}[nosep]
  \item \code{bound} --- the whole crop, $330\times98$ layer pixels.
  \item \code{bound\_trim05} and \code{bound\_trim10} ---
        the crop less a $5\pct$ or $10\pct$ border on each side
        (denominator: the crop's own height and width).
  \item Whole frame --- from the older \code{*\_img.h5} captures.
  \item \code{local\_<patch>} --- the patch grown by $50\pct$ on every side,
        about $51\times51$ layer pixels, fitted at low degree.
  \item \code{sector<k>\_<i>} --- the crop split into $k$ horizontal bands,
        with no continuity enforced between bands.
        The split is along $y$ because that is the axis
        the illumination varies on.
\end{enumerate}

Rungs 4 and 5 turn out to be unusable for criterion P,
for the structural reason given in \S\ref{sec:whyratio}
and reported in \S\ref{sec:localfail}.

Rung 3 remains untested.
The capture suites store only the \code{bound} crop.
Whole-frame raw data exists only in the older \code{*\_img.h5} files,
which carry neither a dark frame nor an LED label.
Testing rung 3 properly needs a configuration switch on the normal
capture path
to archive full-image Bayer layers \emph{and} full-image dark Bayer layers.
That is a change in \code{src/} and has deliberately not been made.

% ============================================================
\section{Results}
\label{sec:results}
% ============================================================

Two capture suites, both of plain colored card filling the stage,
taken 2026-08-03:

\begin{center}
\begin{tabular}{lll}
\toprule
Suite & Stage & Frames scored \\
\midrule
\code{white\_20260803\_1742.h5} & white card stock  & 39 \\
\code{mgray\_20260803\_1736.h5} & gray card stock   & 40 \\
\bottomrule
\end{tabular}
\end{center}

One frame is scored per (LED color $\times$ Bayer layer):
the brightest exposure in the ladder that is not clipped.
\S\ref{sec:ladderselect} discusses that selection rule
and proposes a tightening of it.

\subsection{Repeatability floor}
\label{sec:floor}

Before any residual can be interpreted,
the size of a measurement that means ``zero'' has to be known.
Two comparisons that ought to be flat were measured,
as the $34\times34$ block RMS of a ratio between two normalized frames
(denominator: the frame mean),
over frames brighter than 300\,DN:

\begin{center}
\begin{tabular}{lcc}
\toprule
Comparison & white suite & gray suite \\
\midrule
\code{Gb} against \code{Gr}, same frame   & $0.118\pct$ & $0.101\pct$ \\
Adjacent exposures, same LED              & $0.128\pct$ & $0.173\pct$ \\
\bottomrule
\end{tabular}
\end{center}

Shot noise alone on a 1156-pixel mean predicts about $0.04\pct$,
so these figures are above the pure counting limit
but they are the practical floor of the whole pipeline.
Take $0.12\pct$ as that floor.
Nothing at or below it is a measurement.
The $0.5\pct$ in-sample tolerance sits about $4\times$ above it.

\subsection{Magnitude of the uncorrected non-uniformity}

Across the crop, before correction,
the illumination spans about $+7\pct$ to $-17\pct$ of the crop mean.
The \code{chip2} patch, at the bottom of \code{bound},
sits in the worst of it.
That span is the size of the error the flat field exists to remove.

\subsection{In-sample ladder}

\begin{table}[htbp]
\centering
\caption{White card suite, 39 frames.
Columns as defined in \S\ref{sec:measures}:
\code{P med} and \code{P max} are the median and maximum,
over frames, of each frame's worst-patch $P$;
\code{S med} is the median worst-patch $S$;
\code{S/nz} is that tilt divided by its noise reference;
\code{blk rms} is the block residual.
Terms is the number of free coefficients.}
\label{tab:ladder-white}
\begin{tabular}{llrrrrrrl}
\toprule
Extent & $(d_y,d_x)$ & Terms & P med & P max & S med & S/nz & blk rms & P \\
\midrule
\code{bound} & (1,1) &  3 & 8.86 & 16.35 & 5.67 & 10.62 & 3.77 & fail \\
\code{bound} & (2,1) &  5 & 0.88 &  1.59 & 1.68 &  4.13 & 0.47 & fail \\
\code{bound} & (2,2) &  6 & 0.90 &  1.64 & 1.72 &  4.09 & 0.41 & fail \\
\code{bound} & (3,1) &  7 & \textbf{0.24} & \textbf{0.46} & 1.68 & 4.32 & 0.38 & \textbf{PASS} \\
\code{bound} & (3,2) &  9 & 0.21 &  0.41 & 1.74 &  4.49 & 0.28 & PASS \\
\code{bound} & (3,3) & 10 & 0.21 &  0.41 & 1.76 &  4.60 & 0.28 & PASS \\
\code{bound} & (4,2) & 12 & 0.28 &  0.50 & 1.48 &  3.68 & 0.19 & PASS \\
\code{bound} & (6,2) & 18 & 0.28 &  0.58 & 1.21 &  2.80 & 0.15 & fail \\
\midrule
\code{bound\_trim05} & (2,1) & 5 & 1.11 & 2.50 & 1.91 & 4.73 & 0.45 & fail \\
\code{bound\_trim05} & (3,1) & 7 & 0.39 & 1.04 & 1.79 & 4.51 & 0.39 & fail \\
\code{bound\_trim10} & (3,1) & 7 & 0.93 & 1.97 & 2.15 & 5.57 & 0.49 & fail \\
\midrule
\code{sector3} & (1,1) & 3 & --- & --- & 1.22 & 2.85 & --- & n/a \\
\code{sector3} & (2,1) & 5 & --- & --- & 1.08 & 2.96 & --- & n/a \\
\code{local}   & (1,1) & 3 & --- & --- & 1.11 & 2.59 & --- & n/a \\
\code{local}   & (2,1) & 5 & --- & --- & 1.19 & 3.10 & --- & n/a \\
\code{local}   & (2,2) & 6 & --- & --- & 1.18 & 3.09 & --- & n/a \\
\bottomrule
\end{tabular}
\end{table}

\begin{table}[htbp]
\centering
\caption{Gray card suite, 40 frames.
Same columns.
No rung passes criterion P even in-sample;
see \S\ref{sec:graypuzzle}.}
\label{tab:ladder-gray}
\begin{tabular}{llrrrrrrl}
\toprule
Extent & $(d_y,d_x)$ & Terms & P med & P max & S med & S/nz & blk rms & P \\
\midrule
\code{bound} & (1,1) &  3 & 9.63 & 16.63 & 6.99 & 6.11 & 4.18 & fail \\
\code{bound} & (2,1) &  5 & 1.20 &  7.51 & 3.30 & 2.57 & 0.55 & fail \\
\code{bound} & (2,2) &  6 & 1.12 &  7.47 & 3.30 & 2.57 & 0.49 & fail \\
\code{bound} & (3,1) &  7 & 0.73 &  1.29 & 2.58 & 2.02 & 0.49 & fail \\
\code{bound} & (3,2) &  9 & 0.72 &  1.26 & 2.52 & 1.98 & 0.42 & fail \\
\code{bound} & (3,3) & 10 & 0.72 &  1.26 & 2.73 & 2.12 & 0.42 & fail \\
\code{bound} & (4,2) & 12 & 0.55 &  1.08 & 2.40 & 2.03 & 0.37 & fail \\
\code{bound} & (6,2) & 18 & 0.46 &  1.00 & 2.47 & 2.07 & 0.33 & fail \\
\midrule
\code{bound\_trim05} & (3,1) & 7 & 1.08 & 2.23 & 3.10 & 2.36 & 0.50 & fail \\
\code{bound\_trim10} & (3,1) & 7 & 1.09 & 1.85 & 2.88 & 2.10 & 0.67 & fail \\
\bottomrule
\end{tabular}
\end{table}

Four results follow from these tables.

\textbf{Degree 1 is inadequate; degree 3 is where the improvement stops.}
A plane leaves $8.9$ to $9.6\pct$ median P,
so the surface genuinely curves.
Raising $d_y$ from 2 to 3 gains $0.64\pct$ on the white card.
Raising it from 3 to 6 gains nothing.

\textbf{Raising $d_x$ changes nothing.}
$(3,1)$, $(3,2)$ and $(3,3)$ differ by at most $0.03\pct$ in median P.
This is the direct evidence for the anisotropic basis.

\textbf{Overfitting is visible in-sample.}
On the white card,
P at $(6,2)$ is $0.28\pct$ median and $0.58\pct$ maximum,
which is worse than $(3,1)$'s $0.24$ and $0.46$,
while the block residual continues to fall from $0.38$ to $0.15\pct$.
The added terms are describing the card's own texture,
and they do so near the patches, where it costs most.
A criterion based on pixelwise flatness would have selected $(6,2)$.

\textbf{Trimming the fit extent makes the result worse.}
\code{bound\_trim05} at $(3,1)$ scores $0.39\pct$
against \code{bound}'s $0.24\pct$,
and \code{trim10} scores $0.93\pct$.
The patches sit near the edges of the crop,
so trimming forces the polynomial to extrapolate
exactly where it is being judged.
Rung 2 of the ladder is rejected.

\subsection{Local and sector fits cannot produce a patch ratio}
\label{sec:localfail}

Rungs 4 and 5 report no P at all.
This is the structural consequence derived in \S\ref{sec:whyratio}:
each local or sector field is normalized over its own extent,
so the two normalizing constants do not cancel in the ratio,
and their quotient is exactly the relative brightness of the two extents ---
which is the quantity being measured.

The harness reports these entries as blank
rather than computing a number that would look meaningful
and be arbitrary.
Their S column is still informative,
and local fits do flatten a patch well,
but S is the secondary criterion.
Any scheme using more than one field
needs a step that ties the fields to a common scale
before it can be used at all.

\subsection{Transfer}
\label{sec:transfer}

Fields fitted on the white suite, applied to the gray suite.
Both P columns use the corrected white-patch mean as denominator.
The \code{aligned} columns repeat the transfer
with the field displaced bodily by the whole number of layer rows
that suits it best;
\S\ref{sec:registration} explains why that column exists.

\begin{table}[htbp]
\centering
\caption{Transfer: white $\rightarrow$ gray,
33 (LED $\times$ layer) pairs after exclusions.
Violet excluded, per \S\ref{sec:fluor}.}
\label{tab:transfer}
\begin{tabular}{lrrrrrrrr}
\toprule
 & & \multicolumn{2}{c}{in-sample} & \multicolumn{2}{c}{transfer}
 & \multicolumn{3}{c}{aligned} \\
\cmidrule(lr){3-4}\cmidrule(lr){5-6}\cmidrule(lr){7-9}
$(d_y,d_x)$ & Terms & P med & blk & P med & P max & P med & P max & @row \\
\midrule
(1,1) &  3 & 8.98 & 3.81 & 10.14 & 16.51 & 10.06 & 16.24 & $+14$ \\
(2,1) &  5 & 0.92 & 0.46 & \textbf{0.98} & 1.90 & \textbf{0.79} & 1.66 & $+2$ \\
(3,1) &  7 & 0.23 & 0.38 & 1.34 & 2.05 & 1.13 & 1.95 & $+2$ \\
(4,1) &  9 & 0.27 & 0.32 & 1.41 & 2.15 & 1.24 & 2.01 & $+2$ \\
(4,2) & 12 & 0.27 & 0.19 & 1.40 & 2.14 & 1.24 & 2.00 & $+2$ \\
(5,2) & 15 & 0.26 & 0.16 & 1.32 & 2.07 & 1.13 & 1.90 & $+2$ \\
(6,2) & 18 & 0.28 & 0.15 & 1.38 & 2.20 & 1.20 & 2.03 & $+2$ \\
\bottomrule
\end{tabular}
\end{table}

The difference between the two halves of the table is the point.
In-sample P improves by a factor of four from $(1,1)$ to $(3,1)$
and then stops.
Transfer P is best at $(2,1)$ and grows worse as degree rises.
Criterion T, applied mechanically, selects $(2,1)$.

That selection should not be acted on,
for the reason given in \S\ref{sec:staging}:
the two suites imaged \emph{different physical stages}.
A transfer between captures of two different cards
measures card differences and model error at the same time
and cannot separate them.
Table~\ref{tab:transfer} should be read as an upper bound ---
transfer error is no worse than about $1\pct$ median ---
with the ranking of degrees within it not yet trustworthy.

\subsection{Registration offset}
\label{sec:registration}

The shift scan displaces the fitted field by a whole number of layer rows
and rescores.
At $(2,1)$:

\begin{center}
\begin{tabular}{lrrrrrrrr}
\toprule
Offset (layer rows) & $-8$ & $-4$ & $-2$ & $0$ & $+2$ & $+4$ & $+8$ \\
Median transfer P ($\pct$) & 2.14 & 1.58 & 1.28 & 0.98 & \textbf{0.79}
                           & 0.84 & 1.59 \\
\bottomrule
\end{tabular}
\end{center}

There is a clear minimum at $+2$ layer rows,
which is $+4$ stream pixels,
improving median P from $0.98$ to $0.79\pct$.
That is a registration offset between the two captures ---
the card or the stage moved ---
and no polynomial degree can correct it.
Only re-registering the captures can.

The scan exists because without it the degree comparison is confounded.
A low-degree field can score well on a transfer test
simply by being too smooth to respond to a misalignment.
Comparing degrees at each one's own best alignment removes that advantage.

\subsection{Confounders found and handled}

\subsubsection{Optical brighteners in the white card}
\label{sec:fluor}

Under violet illumination the white card returned $22.4\times$
the signal of the gray card,
where every other LED color agreed on a ratio near $3.7\times$.
(Denominator in both figures: the gray card's dark-subtracted patch mean.)
A sixfold departure from the consensus ratio
is not a reflectance difference.

White card stock contains optical brighteners:
fluorescent dyes that absorb near-ultraviolet light
and re-emit it in the blue.
Under violet LEDs the card is a light \emph{source},
and the surface fitted to it is the emission pattern of the paper,
which will not transfer to any other target.

The harness detects this automatically,
in \code{find\_anomalous\_leds()},
using a threshold of $2\times$ the consensus ratio.
It reports the finding and excludes the offending LED
from the pooled figures.
That is why Table~\ref{tab:transfer} states ``violet excluded''
rather than silently carrying a $20\pct$ outlier.

\subsubsection{Green crosstalk under red LEDs}
\label{sec:crosstalk}

\code{Gb} and \code{Gr} agree to $0.10$--$0.12\pct$ under most LEDs,
which is the repeatability floor.
Under \code{photo\_red}, \code{red} and \code{red\_orange}
they disagree by more than $1\pct$ on the same block-RMS measure
(denominator: the frame mean),
reproducibly at every exposure,
so the disagreement is not noise.

The green photosites that neighbor red pick up different optical and
electrical crosstalk from those that neighbor blue.
The consequence is firm:
fit each Bayer layer separately,
and never average the two greens.
Agreement between \code{Gb} and \code{Gr} remains a useful internal
control, but only for the non-red LEDs.

\subsubsection{Card texture, and how much of the residual it explains}
\label{sec:texture}

The residuals reported above are limited partly by the card,
not only by the model,
and the two can be separated.

Take two exposures of the same LED and layer,
scale each to its own mean,
and form their ratio.
Structure that is fixed on the card appears identically in both frames
and cancels in that ratio;
random noise does not cancel, and adds in quadrature,
so the ratio's spread divided by $\sqrt{2}$ estimates the noise alone.
Subtracting that from a single frame's spread, in quadrature,
leaves the fixed structure.

At the pixel scale, on the sample patch,
with both cards at comparable brightness:

\begin{center}
\begin{tabular}{lrrrr}
\toprule
Suite / LED / layer & mean DN & pixel CV & noise part & fixed part \\
\midrule
white / \code{white} / \code{B}      & 526 & $1.44\pct$ & $1.12\pct$ & $0.89\pct$ \\
white / \code{cyan} / \code{B}       & 632 & $1.34\pct$ & $1.09\pct$ & $0.78\pct$ \\
gray / \code{royal\_blue} / \code{B} & 532 & $4.01\pct$ & $1.29\pct$ & $3.80\pct$ \\
gray / \code{violet} / \code{B}      & 580 & $3.51\pct$ & $1.15\pct$ & $3.32\pct$ \\
\bottomrule
\end{tabular}
\end{center}

(All percentages are of the patch mean.)
At the same signal level the random noise is the same on both cards,
near $1.1\pct$,
as it must be if it is shot noise.
The fixed part is four times larger on the gray card.
The gray card stock is substantially more textured than the white.

At the $34\times34$ patch scale,
after removing the fitted field,
the same separation gives:

\begin{center}
\begin{tabular}{lrrr}
\toprule
Suite / LED / layer & block CV & noise part & fixed part \\
\midrule
white / \code{white} / \code{B}      & $0.394\pct$ & $0.091\pct$ & $0.383\pct$ \\
white / \code{green} / \code{Gr}     & $0.375\pct$ & $0.084\pct$ & $0.365\pct$ \\
white / \code{blue} / \code{Gr}      & $0.353\pct$ & $0.045\pct$ & $0.350\pct$ \\
gray / \code{royal\_blue} / \code{B} & $0.328\pct$ & $0.145\pct$ & $0.294\pct$ \\
\bottomrule
\end{tabular}
\end{center}

(All percentages are of the frame mean, per Eq.~\ref{eq:block}.)
This is the more important of the two tables.
The block residual after a $(3,1)$ fit ---
$0.38\pct$ for the white suite in Table~\ref{tab:ladder-white} ---
is almost entirely fixed structure on the card,
with only $0.05$ to $0.09\pct$ attributable to noise.

Two conclusions follow.
The polynomial is performing better than the block-residual figure
suggests,
because most of that figure is the target rather than the model.
And criterion P, which assumes the imaged card is uniform,
is measuring card structure and fit error together
whenever the card has structure of its own.
Separating them is exactly what the \code{grayR180} capture of
\S\ref{sec:protocol} is for.

\subsubsection{Staging differences between the two suites}
\label{sec:staging}

The white and gray suites are not two captures of one scene.
They used different sheets of card stock, separately placed.
Paper texture, thickness, flatness and placement all differ at once,
and the $+2$-row registration offset confirms that placement did.

There is a second staging issue.
The production chamber has a flat-black stage and curtains
except at the sample,
whereas a card filling the whole stage
scatters substantially more light back into the optical path
than the instrument will normally see.
The surfaces fitted here therefore include stray light
that the instrument does not have.

Neither issue invalidates the analysis.
The in-sample results, the model form, the criteria
and the confounders are all sound.
Both are reasons the transfer figure is an upper bound.
Measuring the true transfer error requires the controlled protocol
of \S\ref{sec:protocol}.

\subsubsection{The gray suite fails P in-sample}
\label{sec:graypuzzle}

Table~\ref{tab:ladder-gray} shows no rung passing criterion P
on the gray card:
best case $0.46\pct$ median but $1.00\pct$ maximum at $(6,2)$,
and $0.73/1.29\pct$ at $(3,1)$.
The white card passes comfortably at the same degree.

Signal level is the likeliest explanation.
The gray card returns $0.27\times$ the white card's signal,
so its dim layers sit where noise is a much larger fraction
of the patch mean.
\S\ref{sec:ladderselect} tabulates this,
and the worst transfer pairs in \S\ref{sec:transfer}
are readings at 7 to 14\,DN on the gray card.

Card structure is the other candidate,
but \S\ref{sec:texture} weighs against it:
at patch scale the gray card's fixed structure
($0.29\pct$) is no worse than the white card's ($0.35$ to $0.38\pct$),
even though at pixel scale it is four times larger.
The texture appears to be fine-grained enough
to average out over a $34\times34$ patch.

The \code{grayA}/\code{grayA2} pair of \S\ref{sec:protocol}
resolves this,
and the case for re-metering the exposures against the calibration card
is strong either way.

\subsubsection{Enlarging \code{bound} makes the result worse}

This was tested directly,
because enlarging \code{bound} had been proposed
as a way to give the polynomial more support.

At fixed degree, patch error grows as the fit area grows.
On the gray card the worst-patch P went from $0.46\pct$
over a 330-row crop to $2.61\pct$ over an enlarged one;
on a light-gray capture, from $0.17\pct$ to $3.06\pct$.
Raising the degree did not recover it.
The result was confirmed independently on a later 530-row capture:
fitting over 530 rows costs $1.3$ to $2.1\times$ more residual
within the same central region
than fitting over 330 rows does.

A polynomial distributes its degrees of freedom
over the whole fit extent.
Enlarging the extent asks the same 7 terms to describe more of a curve,
and the added territory is territory nothing is measured in.
The recommendation is the opposite one:
move the reference patches inward, toward the sample,
rather than moving \code{bound} outward.
\code{chip2} sits at the bottom of \code{bound},
where illumination has fallen $17\pct$,
and it is the worst patch in every transfer test.

\subsection{Exposure ladder selection}
\label{sec:ladderselect}

\code{capture\_runner} records several exposures per LED in a suite,
and some are expected to saturate somewhere in the frame.
Discarding those and requiring only that one usable exposure survive
is the right policy,
with one refinement.

\textbf{The selection must be per (LED, layer), not per LED.}
The harness already does this,
in \code{select\_frames()}:
for each (LED color, Bayer layer) independently
it takes the unclipped exposure with the highest mean.
Selecting per LED would be wrong,
because the four layers of one frame differ enormously in level.
Under the green LED at exposure \code{e2},
\code{Gr} and \code{Gb} are $96\pct$ masked
while \code{R} sits at 66\,DN and \code{B} at 220\,DN,
both perfectly usable.
Discarding the whole frame would throw away good data
in two layers out of four.

\textbf{``Unclipped'' is necessary but not sufficient.}
An unclipped frame can still be too dim to be worth fitting.
Averaging $n$ pixels reduces relative noise by $\sqrt{n}$,
so the relative noise on a patch mean is
\begin{equation}
  \varepsilon_i \;=\; \frac{\sigma_i}{\sqrt{n_i}} ,
  \label{eq:patchnoise}
\end{equation}
where $\sigma_i$ is the pixel coefficient of variation within the patch,
in percent, and $n_i = 1156$.
Measured on the sample patch of the selected frames:

\begin{center}
\begin{tabular}{lrrr}
\toprule
Suite & patch mean, DN & pixel CV $\sigma$ & patch-mean noise $\varepsilon$ \\
\midrule
white & 27.5  & $5.87\pct$ & $0.173\pct$ \\
white & 70.9  & $3.20\pct$ & $0.094\pct$ \\
white & 241.3 & $2.00\pct$ & $0.059\pct$ \\
white & 903.5 & $1.12\pct$ & $0.033\pct$ \\
\midrule
gray  & 7.8   & $15.59\pct$ & $0.459\pct$ \\
gray  & 45.2  & $5.66\pct$  & $0.167\pct$ \\
gray  & 310.8 & $4.46\pct$  & $0.131\pct$ \\
gray  & 579.8 & $3.51\pct$  & $0.103\pct$ \\
\bottomrule
\end{tabular}
\end{center}

(All percentages are of the sample patch's own mean.)
The gray card never reaches $\varepsilon = 0.1\pct$ at any exposure
in the suite,
while the white card reaches it above about 70\,DN.
That is a concrete, quantitative statement of why the gray suite
underperforms,
and it is a direct argument for re-metering the exposure ladder
against the calibration card.

\textbf{Proposed success criterion.}
A suite is adequate for flat-field work when,
for every (LED color, Bayer layer) combination the analysis consumes,
at least one exposure in the ladder satisfies both of:

\begin{enumerate}[nosep]
  \item masked fraction $\le 1\pct$ of the fit extent
        (already enforced, \S\ref{sec:masking});
  \item $\varepsilon \le 0.1\pct$ by Eq.~\ref{eq:patchnoise},
        which at $n = 1156$ means a pixel CV of $\sigma \le 3.4\pct$
        (not yet enforced).
\end{enumerate}

Exposures that fail either test are discarded, not repaired,
and there is no requirement that one exposure serve all four layers.
The second test is a small addition to \code{select\_frames()}
and would turn the present silent underperformance
into an explicit report.

\textbf{Which combinations must be consumed.}
The five worst transfer pairs of \S\ref{sec:transfer} ---
\code{amber/B}, \code{blue/R}, \code{photo\_red/B},
\code{royal\_blue/R}, \code{red\_orange/B} ---
are all combinations in which the LED color and the Bayer layer
do not match,
and all of them read between 7 and 14\,DN on the destination card.
They are crosstalk readings carrying almost no signal.
If the measurement consumes only the matching combinations,
the success criterion need only apply to those,
and the pooled transfer figures should improve substantially
once the mismatched ones are excluded.
Testing that requires a layer filter on the harness,
which it does not yet have.

% ============================================================
\section{Choices made}
\label{sec:choices}
% ============================================================

\begin{center}
\small
\begin{tabular}{lll}
\toprule
Question & Choice & Basis \\
\midrule
Correction form   & divide by unit-mean gain & \S\ref{sec:divide} \\
Basis             & anisotropic $(d_y,d_x)$  & row vs column variation, $25$ vs $2\pct$ \\
Default degree    & $(3,1)$, 7 terms         & Table~\ref{tab:ladder-white} \\
Fit extent        & the whole crop           & trims and enlargements both worse \\
Granularity       & per LED color, per layer & $2.1\pct$ and $0.96\pct$ shape differences \\
Normalization     & mean over the fit extent & \S\ref{sec:whyratio} \\
Coordinates       & $[-1,1]$ over \code{bound} always & comparable coefficients \\
Degree selection  & aligned transfer P       & criterion T \\
Runtime form      & per-patch scalar gains   & \S\ref{sec:scalar} \\
Frame selection   & best unclipped, per (LED, layer) & \S\ref{sec:ladderselect} \\
Calibration card  & brightener-free and smooth & \S\ref{sec:fluor}, \S\ref{sec:texture} \\
Storage           & coefficients authoritative & maps and scalars derived \\
\bottomrule
\end{tabular}
\end{center}

$(3,1)$ is the default, not a settled answer.
Criterion T on the present data prefers $(2,1)$,
but on data that cannot separate model error from card differences.
The choice between $(2,1)$ and $(3,1)$
is to be re-run on the controlled suites.

% ============================================================
\section{Implementation}
\label{sec:implementation}
% ============================================================

\subsection{\code{devdesk/flatfield.py}}

A desktop-only evaluation harness,
sibling of \code{hdf5inspect.py} and \code{roi\_darkdiff.py}
in the same directory,
sharing their conventions:
argparse mode flags,
\code{Hdf5Reader} from \code{src/hdf5read.py},
bare filenames resolved against \code{captures/},
results into \code{devdesk/devdesk.sqlite}.
It makes no changes to \code{src/} behavior.

Principal types:

\begin{description}[style=nextline, leftmargin=1em]
\item[\code{Frame}]
  One dark-subtracted, masked layer,
  with its LED, exposure label, exposure in microseconds and dark link.
  \code{.usable} is false when more than $1\pct$ of pixels are masked.
\item[\code{Patches}]
  The active patch rectangles in layer coordinates,
  plus which one holds the white-reference role.
  Loaded through \code{roi.load\_active()} and
  \code{roi.to\_layer\_rect\_within()};
  never hardcoded.
\item[\code{FlatField}]
  Coefficients, degrees, normalization domain, fit extent,
  the window it covers, and \code{norm\_mean} $= \mean{F}{\Omega}$.
  Provides \code{gain()}, \code{correct()}, \code{gain\_for\_patch()}
  and \code{coef\_json()}.
\item[\code{Model}]
  A named collection of one or more \code{FlatField}s
  with \code{field\_for(rect)},
  so that single-field and multi-field extents
  are scored through one interface.
\item[\code{Score}]
  P, S, $S/S^{\mathrm{nz}}$, block residual,
  and the count of patches no field could cover without extrapolating.
\end{description}

The \code{covers} field on \code{FlatField} deserves a note,
because getting it wrong produced a wrong answer once.
\code{bound} extents cover the whole crop;
\code{local\_*} and \code{sector*} extents cover only their own window.
Without that distinction,
trimmed fits were scored on whichever patches happened to fall
inside the trim window,
which is a smaller and easier subset,
and trimming looked better than it is.

\subsection{Storage}

Coefficients are the source of truth;
maps and scalars are regenerable caches.

Table \code{flatfield} in \code{devdesk/devdesk.sqlite},
keyed on (\code{h5file}, \code{led}, \code{eno}, \code{bayerlayer},
\code{extent}, \code{deg\_y}, \code{deg\_x}),
holds \code{coef\_json},
the normalization domain
(\code{domain\_rows}, \code{domain\_cols}, \code{bound\_*}, \code{frame\_*}),
\code{norm\_mean},
and the fit statistics
(\code{n\_fit}, \code{n\_masked}, \code{fit\_rms\_pct},
\code{block\_rms\_pct}, \code{patch\_dev\_max\_pct}).

Derived gain maps go to

\begin{footnotesize}
\begin{verbatim}
    devdesk/flatfield/<suite_id>/<led>_<ch>_<roi>_gain.npy
\end{verbatim}
\end{footnotesize}

\noindent
holding the field evaluated over one patch rectangle,
divided by its own mean.
This is the artifact \code{colorimeter\_methods.tex} called
\code{color\_bayer\_rect\_trend.npy},
changed from additive to multiplicative,
and regenerable from the coefficients whenever a patch moves.
That regenerability is the reason for keeping the coefficients
authoritative.

\S\ref{sec:scalar} adds a third, smaller artifact for the runtime path:
the per-patch scalar $\bar{g}_i$,
which is the mean of that same gain map.

\subsection{Archive metadata (capture schema 2)}

The analysis needs to know the \code{bound} rectangle
a capture was taken with.
The stored layer shape gives its width and height;
it cannot give its $x$ and $y$,
and a uniform card carries no features to register against.

Consulting the local \code{roi} table is not a substitute.
The analysis machine's \code{colorimeter.sqlite} is a copy,
and an ROI edit made on the instrument never reaches it.
This produced a real false alarm:
an apparent fault in the HDF5 writer
that was in fact a stale desktop database.

Archives now describe themselves.
\code{capture\_h5.SCHEMA\_VERSION} is 2,
and each archive carries:

\begin{itemize}[nosep]
  \item root attributes \code{bound\_x}, \code{bound\_y},
        \code{bound\_w}, \code{bound\_h},
        which are cheap to read and sufficient to interpret the layers;
  \item a \code{/meta} JSON record with the full patch set,
        camera settings,
        and rows from the \code{device} and \code{samples} tables.
\end{itemize}

\code{src/capture\_meta.py} collects that record declaratively.
Column lists and field extractors are tuples at the top of the module,
so adding a field is a one-line edit in one place.
\code{collect()} never raises;
a block that fails becomes \code{\{"error": \ldots\}},
because a metadata problem must never cost a capture.

\code{flatfield.py} reads \code{archive\_bound()} first
and falls back to the \code{roi} table only for schema-1 archives,
printing a warning when it does.
A \code{--bound X,Y,W,H} override handles archives
taken against geometry this database does not have.

\medskip
\noindent
\textbf{Deployment note.}
The schema-2 writer must be deployed to the Pi
before the protocol captures of \S\ref{sec:protocol} are taken.
Archives written by the current Pi code are still schema 1
and will not carry their own geometry.

% ============================================================
\section{Image sampling protocol}
\label{sec:protocol}
% ============================================================

The binding uncertainty is the transfer error,
and it cannot be resolved from the existing captures,
because those used two different stages.
Four short suites on one card resolve it.

\subsection{Card selection}

The card must satisfy two independent requirements,
and the earlier recommendation to use gray stock addressed only the first.

\begin{enumerate}[nosep]
  \item \textbf{Free of optical brighteners.}
        White stock fluoresces under violet (\S\ref{sec:fluor}),
        and its surface under that LED is the paper's emission
        rather than the instrument's illumination.
        Check by confirming that the violet LED's brightness ratio
        against a reference card
        sits with the other LEDs rather than several times above them.
  \item \textbf{Smooth at the pixel scale.}
        The gray stock measured in \S\ref{sec:texture}
        carries $3.3$ to $3.8\pct$ fixed texture at pixel scale,
        four times the white stock's.
        Check by the two-exposure ratio method of \S\ref{sec:texture},
        which separates fixed texture from random noise;
        aim for a fixed part below $1.5\pct$ of the patch mean.
\end{enumerate}

The gray stock's texture largely averages out over a $34\times34$ patch,
so it is usable,
but a card meeting both requirements gives a cleaner measurement
of the instrument.

\subsection{Preparation}

\begin{enumerate}[nosep]
  \item Select and verify the card as above.
  \item \textbf{Mark one corner} of the card,
        indelibly and visibly in frame if possible.
        The \code{grayR180} suite depends on knowing the orientation.
  \item \textbf{Re-meter the exposure ladder against this card},
        so that the maximum DN anywhere inside \code{bound}
        is $\le 850$\,DN in every layer.
        Metering on a center patch is what oversaturated
        the earlier white suite.
        A ceiling of 850 leaves margin below the 1000\,DN mask threshold
        for the brighter corners.
  \item Confirm that the ladder meets the success criterion of
        \S\ref{sec:ladderselect}:
        for every (LED, layer) combination consumed,
        at least one exposure with $\sigma \le 3.4\pct$
        on the sample patch.
  \item \textbf{Deploy schema-2 capture code to the Pi} first,
        so that each archive records its own \code{bound}.
  \item Let the LEDs and the enclosure reach thermal steady state
        before the first suite,
        and take all four suites in one sitting.
\end{enumerate}

\subsection{The four suites}

\begin{center}
\begin{tabular}{lll}
\toprule
Suite & Action before capture & What it isolates \\
\midrule
\code{grayA}    & place the card, capture & baseline \\
\code{grayA2}   & \textbf{touch nothing}, capture again & noise floor alone \\
\code{grayB}    & lift the card, replace it carefully & placement repeatability \\
\code{grayR180} & rotate the same card $180^{\circ}$ in plane
                & paper against instrument \\
\bottomrule
\end{tabular}
\end{center}

Each suite is the full LED set with the re-metered ladder.
It is the same capture the instrument already performs,
with a card on the stage.

\subsection{What each suite answers}

\code{grayA} $\rightarrow$ \code{grayA2} should transfer
at the repeatability floor, about $0.12\pct$ of the white patch mean.
If it does not,
something in the pipeline is contributing error of its own
and every other figure here is suspect.
This is the control.

\code{grayA} $\rightarrow$ \code{grayB} adds placement.
The difference between this result and \code{grayA2}'s
is the cost of re-seating a card,
and it sets a floor on how well any calibration
can be verified by re-imaging a target.

\code{grayA} $\rightarrow$ \code{grayR180} is the decisive one.
Rotating the card $180^{\circ}$ leaves the instrument's illumination
pattern untouched
but rotates the paper's own structure with the card.
Therefore:

\begin{itemize}[nosep]
  \item Structure that stays in place under rotation
        belongs to the instrument.
        It is what the polynomial is for,
        and it is correctable.
  \item Structure that rotates with the card belongs to the paper.
        No flat field can remove it,
        and it is a hard floor on how well criterion P
        can ever be verified,
        because P assumes the card is uniform
        and this measures the extent to which it is not.
\end{itemize}

That decomposition converts Table~\ref{tab:transfer}'s upper bound
into an actual measurement of model error,
and it settles the choice between $(2,1)$ and $(3,1)$.
\S\ref{sec:texture} has already shown the effect is large enough
at the pixel scale to be worth separating,
and this is the measurement that separates it at the patch scale.

\subsection{Optional additions}

If a configuration switch for full-image Bayer layers
plus full-image dark layers is added to the normal capture path,
take one extra \code{grayA} suite with it enabled.
That makes ladder rung 3 testable
and answers directly whether \code{bound}
is a large enough support region.

% ============================================================
\section{Running the programs}
\label{sec:running}
% ============================================================

Use the project virtual environment.
The system Python lacks \code{h5py}, \code{cv2} and \code{matplotlib}.

\begin{verbatim}
cd /home/bill/Projects/Colorimeter
.venv/bin/python3 devdesk/flatfield.py --selftest
\end{verbatim}

\noindent
Modes:

\begin{center}
\begin{tabular}{ll}
\toprule
Flag & Purpose \\
\midrule
\code{-f}, \code{--fit}      & fit the default $(3,1)$ over \code{bound};
                               print coefficients \\
\code{-c}, \code{--compare}  & the extent $\times$ degree ladder, in-sample \\
\code{-t}, \code{--transfer} & fit on FILE, apply to FILE2 \\
\code{-p}, \code{--plot}     & profiles, residual heatmap,
                               error-against-degree curves \\
\code{-w}, \code{--write}    & store coefficients and derived gain maps \\
\bottomrule
\end{tabular}
\end{center}

\noindent
Filters and overrides:
\code{--led COLOR} (repeatable),
\code{--exp LABEL} (repeatable),
\code{--deg-y}, \code{--deg-x},
\code{--bound X,Y,W,H},
\code{--captures-dir}, \code{--db}, \code{--gain-dir},
and \code{-o FILE.png} to save a \code{-p} figure instead of showing it.

\subsection{Fitting one LED}

\begin{verbatim}
.venv/bin/python3 devdesk/flatfield.py white_20260803_1742.h5 -f --led green
\end{verbatim}

\begin{footnotesize}
\begin{verbatim}
white_20260803_1742.h5: flat-field fits, (deg_y, deg_x) = (3, 1), 7 terms
  basis: y^0x^0 y^0x^1 y^1x^0 y^1x^1 y^2x^0 y^2x^1 y^3x^0
  led      e#  ch   exp_us    mean  masked%   span%  resid%  coef[0..2]
  green    e0  R      5291    29.1    0.00%   29.1%   5.86%    30.47  -0.43   -3.64
  green    e0  Gr     5291   514.1    0.00%   31.2%   1.59%   538.65  -6.52  -69.48
  green    e0  Gb     5291   514.0    0.00%   31.3%   1.59%   538.66  -6.62  -69.87
  green    e0  B      5291    97.3    0.00%   35.8%   3.02%   102.40  -1.51  -16.22
  green    e1  Gr     7937   770.9    0.00%   31.1%   1.45%   808.37  -9.46 -102.98
  green    e2  Gr    11906   918.1   96.24%  SKIPPED - over 1% masked
\end{verbatim}
\end{footnotesize}

\noindent
Column meanings:

\begin{description}[style=nextline, leftmargin=1em, nosep]
\item[\code{mean}]
  $\mean{S}{\mathrm{crop}}$, the dark-subtracted mean over the crop, in DN.
  Use it to judge whether a row is worth reading at all.
  At 29\,DN the \code{R} layer under a green LED is nearly all noise,
  which is why its \code{resid\%} is $5.86$
  while the well-exposed \code{Gr} row's is $1.59$.
\item[\code{masked\%}]
  Percent of crop pixels excluded.
  Denominator: all pixels in the crop.
  Above $1\pct$ the row is \code{SKIPPED}.
\item[\code{span\%}]
  $(\max F - \min F)$ divided by $\mean{F}{\Omega}$,
  over the fitted surface.
  It is the size of the illumination variation the fit describes,
  about $30\pct$ here.
\item[\code{resid\%}]
  Root-mean-square of the pixelwise fit residual,
  as a percentage of the frame mean.
  This is a \emph{pixelwise} figure and is dominated by shot noise.
  Do not compare it with criterion P,
  which is a patch-mean figure and roughly $\sqrt{1156}$ times smaller
  for the same underlying error.
\item[\code{coef[0..2]}]
  The first three coefficients:
  the constant, the $x$ term and the $y$ term.
  The $y$ coefficient is about $10\times$ the $x$ coefficient
  with the same sign pattern in every row,
  which is the anisotropy in raw form.
\end{description}

\subsection{The ladder}

\begin{verbatim}
.venv/bin/python3 devdesk/flatfield.py mgray_20260803_1736.h5 -c
\end{verbatim}

Prints the repeatability floor,
the patch positions in use,
then the content of Table~\ref{tab:ladder-gray}.
Read the floor block first;
it states what the smallest meaningful number in the table is.

\subsection{Transfer}

\begin{verbatim}
.venv/bin/python3 devdesk/flatfield.py grayA_<stamp>.h5 grayR180_<stamp>.h5 -t
\end{verbatim}

Prints the fluorescence exclusion list,
the degree table (Table~\ref{tab:transfer}),
the criterion-T recommendation,
the shift scan,
and the five worst (LED $\times$ layer) pairs
with their source and destination brightness in DN.

\subsection{Writing results}

\begin{verbatim}
.venv/bin/python3 devdesk/flatfield.py grayA_<stamp>.h5 -w
\end{verbatim}

Writes \code{flatfield} rows and \code{*\_gain.npy} maps,
then verifies the round trip
by rebuilding a gain map from the stored \code{coef\_json}
and comparing it with the file.

% ============================================================
\section{Interpreting results}
\label{sec:interpreting}
% ============================================================

\subsection{Reading order}

\begin{enumerate}[nosep]
  \item \textbf{The floor block.}
        If \code{Gb} against \code{Gr} is not near $0.12\pct$,
        stop.
        The pipeline is contributing error
        and nothing below it is trustworthy.
  \item \textbf{The exclusion warnings.}
        A fluorescence flag or a \code{SKIPPED} row
        changes which frames the pooled figures cover.
  \item \textbf{The patch list.}
        Confirm the patch positions are the intended ones;
        a stale \code{roi} table moves them without any other symptom.
  \item \textbf{Transfer P, not in-sample P},
        for any decision about degree.
  \item \textbf{\code{P max}, not \code{P med}},
        for any decision about acceptance.
        The median conceals the one patch that fails.
\end{enumerate}

\subsection{Decision rules}

\begin{description}[style=nextline, leftmargin=1em]

\item[Transfer P falls with degree, then flattens.]
  The expected pattern.
  Take the degree at which it flattens, per criterion T.

\item[Transfer P rises with degree while in-sample P falls.]
  Overfitting.
  Take the lower degree.
  This is what Table~\ref{tab:transfer} shows above $(2,1)$.

\item[The shift scan has a minimum away from zero.]
  The two captures are not registered to each other.
  Correct the staging or the registration;
  changing the polynomial cannot help.
  Compare degrees on the \code{aligned} columns only.

\item[S is large but \code{S/nz} is near 1.]
  The tilt is a plane fitted through noise, not a real tilt.
  Ignore it.
  This is common in dim layers,
  where the pixel coefficient of variation reaches $6\pct$.

\item[S is large and \code{S/nz} is well above 1.]
  A real residual tilt that no whole-crop polynomial removed.
  It does not bias the patch mean,
  so P can still pass,
  but it will inflate Layer 3's uniformity flag.

\item[Block residual keeps improving while P stalls or worsens.]
  The added terms are describing texture, not illumination.
  Trust P.
  Compare the block residual against \S\ref{sec:texture}'s figures
  before concluding the model is at fault:
  most of that residual is the card.

\item[P is blank.]
  A multi-field model.
  Not a missing number; see \S\ref{sec:localfail}.

\item[A single LED is a large outlier.]
  Check fluorescence first, then signal level.
  The worst-pairs listing prints source and destination DN
  precisely so that a $1.6\pct$ error at 11\,DN
  can be recognized as noise rather than field error.

\end{description}

\subsection{What adequate will look like}

On the controlled suites, the flat field is adequate when all of:

\begin{itemize}[nosep]
  \item \code{grayA} $\rightarrow$ \code{grayA2} transfer \code{P max}
        $\le 0.3\pct$ of the corrected white patch mean,
        which is near the floor;
  \item \code{grayA} $\rightarrow$ \code{grayB} transfer \code{P max}
        $\le 1.0\pct$;
  \item in \code{grayA} $\rightarrow$ \code{grayR180},
        the part of the residual that stays fixed under rotation
        is $\le 1.0\pct$,
        the rotating part being card structure rather than model error;
  \item the shift scan on \code{grayA} $\rightarrow$ \code{grayA2}
        has its minimum at zero rows;
  \item the chosen degree is the point at which aligned transfer P
        flattens,
        not the point at which in-sample residual flattens.
\end{itemize}

% ============================================================
\section{How this fits into next steps}
\label{sec:next}
% ============================================================

\subsection{Settled}

\begin{itemize}[nosep]
  \item Model form:
        multiplicative correction,
        anisotropic polynomial,
        one per LED color and per Bayer layer,
        fitted over the whole crop,
        normalized over its fit extent.
  \item Criteria defined and implemented:
        P, S with its noise reference, and T on aligned transfer.
  \item The correction reduces at measurement time
        to two scalars per patch pair (\S\ref{sec:scalar}),
        with the error of that reduction bounded
        by a product of within-patch variations.
  \item Extent ladder mostly resolved:
        the whole crop wins;
        trims are worse;
        multi-field schemes cannot produce a patch ratio
        without a cross-normalization step;
        whole-frame is untestable with current archives.
  \item Confounders identified and handled in code:
        saturation, optical brighteners, red-LED green crosstalk,
        registration offset,
        and card texture.
  \item Archives made self-describing (capture schema 2),
        so analysis no longer depends on a database copy
        staying in step.
  \item Patch selection converted from hardcoded names to roles,
        so that renaming a patch cannot silently misdirect the analysis.
\end{itemize}

\subsection{Immediately next}

\begin{enumerate}[nosep]
  \item Deploy schema-2 capture code to the Pi.
  \item Run \code{src/arbiter.py}'s self-test on the Pi.
        It cannot run on the desktop, which lacks \code{gpiozero},
        so the role-based edits to it are compile-checked
        but not executed.
  \item Add the $\varepsilon \le 0.1\pct$ test to
        \code{select\_frames()} (\S\ref{sec:ladderselect}),
        and a layer filter so that mismatched
        (LED, layer) combinations can be excluded from pooled figures.
  \item Select and verify a calibration card
        against both requirements of \S\ref{sec:protocol}.
  \item Re-meter the exposure ladder against it to the 850\,DN ceiling.
  \item Capture \code{grayA}, \code{grayA2}, \code{grayB},
        \code{grayR180}.
  \item Re-run \code{-c} and \code{-t} across them,
        and settle $(2,1)$ against $(3,1)$
        on data that separates model error from card structure.
\end{enumerate}

\subsection{Then}

\begin{itemize}[nosep]
  \item \textbf{Promote the evaluated model into \code{src/}.}
        The harness stays in \code{devdesk/} as the evaluation tool.
        The instrument needs only to look up two stored scalars
        per patch pair,
        with the map generation available for diagnostics.
  \item \textbf{Measure within-patch reflectance variation on developed
        strips},
        which is the one unmeasured factor
        in the bound of \S\ref{sec:scalar}.
  \item \textbf{Decide the calibration cadence:}
        how often the field is refitted, and against what ---
        a stored card, or the white reference patch on each strip.
  \item \textbf{Move the reference patches inward.}
        \code{chip2} sits where illumination has fallen $17\pct$,
        it has eight times the sample patch's within-patch gain variation,
        and it is the worst patch in every test.
        Bringing the references nearer the sample
        reduces the flat field's workload substantially.
        This is a layout change with consequences beyond this task.
  \item \textbf{Bayesian update within the patch:}
        a stored field as prior,
        updated by the few measurements each run actually makes,
        rather than re-estimated from scratch.
        Deferred deliberately,
        because it is worth building only once the prior's quality
        is known,
        and that is what the protocol captures establish.
  \item \textbf{Full-frame capture switch}, to close ladder rung 3.
\end{itemize}

\subsection{Relation to the other documents}

This is Layer 1 of \code{colorimeter\_plan.tex},
with two amendments to \code{colorimeter\_methods.tex}:
the correction divides rather than subtracts (\S\ref{sec:divide}),
and the stored artifact is a multiplicative gain
derived from authoritative coefficients
rather than an additive trend array (\S\ref{sec:implementation}).
Layer 3's strip-uniformity flag is the consumer of criterion S.

\end{document}
